% ============================================================
% Section 3: Problem Formulation and Preliminaries   (draft v1, 23 Oct)
% Packages used: amsmath, amssymb, booktabs (all loaded by acmart).
% Notation: user embedding \mathbf{p}_u, item embedding \mathbf{q}_i,
%           log-popularity \ell_i (NOT p_i; p is reserved for users).
% ============================================================
\section{Problem Formulation and Preliminaries}
\label{sec:prelim}

We fix the recommendation setting (\S\ref{sec:prelim-setting}) and formalise embedding sharing as a fixed many-to-one map onto a smaller table (\S\ref{sec:prelim-sharing}). We then define the two training objectives and four scoring rules we compare (\S\ref{sec:prelim-objectives}). Finally, we state the problem together with the quantities used to diagnose it (\S\ref{sec:prelim-problem}).

\subsection{Recommendation Setting and Notation}
\label{sec:prelim-setting}

Let $\mathcal{U}$ and $\mathcal{I}$ be the sets of users and items, with $n_U=|\mathcal{U}|$ and $n_I=|\mathcal{I}|$. Implicit feedback is a set of observed interactions $\mathcal{E}\subseteq\mathcal{U}\times\mathcal{I}$. It is split into training, validation and test sets, $\mathcal{E}_{\mathrm{tr}}$, $\mathcal{E}_{\mathrm{va}}$ and $\mathcal{E}_{\mathrm{te}}$. The \emph{popularity} of item $i$ is its training count, and we write $\ell_i$ for its log-popularity:
\begin{equation}
  c_i = \bigl|\{u : (u,i)\in\mathcal{E}_{\mathrm{tr}}\}\bigr|,
  \qquad
  \ell_i = \log(1+c_i).
  \label{eq:popularity}
\end{equation}
Only training interactions enter $c_i$. A model assigns each user an embedding $\mathbf{p}_u\in\mathbb{R}^d$ and each item an embedding $\mathbf{q}_i\in\mathbb{R}^d$, and scores a pair with a rule $s(u,i)$ (\S\ref{sec:prelim-objectives}). It recommends the $K$ highest-scoring items that $u$ has not interacted with in training; at test time, validation items are also excluded. We use matrix factorisation (MF), in which $\mathbf{p}_u$ and $\mathbf{q}_i$ are read directly from embedding tables. We also use LightGCN~\cite{LightGCN}, in which $\mathbf{p}_u$ and $\mathbf{q}_i$ are the outputs of $L$ rounds of propagation over the normalised training graph. All later definitions apply to these final vectors unchanged.

\subsection{Embedding Sharing Under a Memory Budget}
\label{sec:prelim-sharing}

A \emph{full} model stores one row per ID, in tables $\mathbf{E}_U\in\mathbb{R}^{n_U\times d}$ and $\mathbf{E}_I\in\mathbb{R}^{n_I\times d}$, at a cost of $(n_U+n_I)\,d$ parameters. Embedding sharing replaces these with smaller tables $\mathbf{E}_U\in\mathbb{R}^{B_U\times d}$ and $\mathbf{E}_I\in\mathbb{R}^{B_I\times d}$, where $B_U\ll n_U$ and $B_I\ll n_I$. It also fixes, before training, the assignment maps $h_U:\mathcal{U}\to[B_U]$ and $h_I:\mathcal{I}\to[B_I]$. Under \emph{single-row sharing},
\begin{equation}
  \mathbf{p}_u = \mathbf{E}_U[h_U(u)],
  \qquad
  \mathbf{q}_i = \mathbf{E}_I[h_I(i)],
  \label{eq:sharing}
\end{equation}
and the table costs $(B_U+B_I)\,d$ parameters, independent of the catalog size.

The \emph{bucket} of item $i$ is $b(i)=h_I(i)$. The items in bucket $b$ are $\mathcal{I}_b=\{i : b(i)=b\}$. The bucket's \emph{traffic} is the sum of its items' counts,
\begin{equation}
  C_b=\sum_{i\in\mathcal{I}_b} c_i .
  \label{eq:traffic}
\end{equation}
\emph{Two-row} schemes give each item two codes and sum the two rows, $\mathbf{q}_i=\mathbf{E}_I[h^{(1)}_I(i)]+\mathbf{E}_I[h^{(2)}_I(i)]$, and likewise for users.

We study four assignment schemes.
\begin{itemize}
  \item \textbf{GraphHash}~\cite{GraphHash} partitions the bipartite training graph by Louvain modularity maximisation, with resolution 200. Each user and each item is assigned the row of its cluster. This is our main setting.
  \item \textbf{Random} hashing assigns IDs uniformly at random, using GraphHash's bucket counts.
  \item \textbf{DoubleFreq} uses frequency-based double hashing~\cite{Zhang2020freq}, a two-row scheme in which popular IDs receive near-unique code pairs.
  \item \textbf{DoubleGraph} is a two-row scheme that sums a universal-hash code and a Louvain-cluster code. We use it as a structure-aware substitute for BACO~\cite{BACO}.
\end{itemize}

On Gowalla, Yelp2018 and AmazonBook, GraphHash yields 5{,}806, 7{,}656 and 9{,}804 item buckets, with 7.1, 5.0 and 9.3 items per bucket on average. The shared table is 16.4\%, 21.9\% and 13.6\% of the size of the full table.

Single-row sharing has an immediate consequence: whenever $b(i)=b(j)$, we have $\mathbf{q}_i=\mathbf{q}_j$. Any item-level quantity the model computes from its item embedding, in particular the magnitude $\|\mathbf{q}_i\|$, is therefore constant within a bucket.

\subsection{Training Objectives and Scoring Rules}
\label{sec:prelim-objectives}

We compare two objectives that differ in whether they see embedding magnitude. Let $\mathcal{B}\subset\mathcal{E}_{\mathrm{tr}}$ be a batch.

\emph{BPR}~\cite{BPR} draws one negative item $j$ uniformly at random for each positive pair and minimises
\begin{equation}
  \mathcal{L}_{\mathrm{BPR}}
  = -\frac{1}{|\mathcal{B}|}\sum_{(u,i)\in\mathcal{B}}
    \log\sigma\!\bigl(\langle\mathbf{p}_u,\mathbf{q}_i\rangle-\langle\mathbf{p}_u,\mathbf{q}_j\rangle\bigr).
  \label{eq:bpr}
\end{equation}

\emph{DirectAU}~\cite{DirectAU} acts on normalised embeddings $\tilde{\mathbf{x}}=\mathbf{x}/\|\mathbf{x}\|$. It combines an alignment term with an in-batch uniformity term:
\begin{equation}
\begin{aligned}
  \mathcal{L}_{\mathrm{DAU}}
  = \frac{1}{|\mathcal{B}|}\sum_{(u,i)\in\mathcal{B}}
      \bigl\|\tilde{\mathbf{p}}_u-\tilde{\mathbf{q}}_i\bigr\|^2
  + \frac{\gamma}{2}\Bigl[
      &\log \mathbb{E}_{u,u'}\, e^{-2\|\tilde{\mathbf{p}}_u-\tilde{\mathbf{p}}_{u'}\|^2} \\
      +\,&\log \mathbb{E}_{i,i'}\, e^{-2\|\tilde{\mathbf{q}}_i-\tilde{\mathbf{q}}_{i'}\|^2}
    \Bigr],
\end{aligned}
\label{eq:dau}
\end{equation}
with $\gamma=1$ and expectations taken over distinct pairs in $\mathcal{B}$.

The distinction that matters here is that $\mathcal{L}_{\mathrm{BPR}}$ depends on the magnitudes $\|\mathbf{p}_u\|$ and $\|\mathbf{q}_i\|$, whereas $\mathcal{L}_{\mathrm{DAU}}$ depends on the embeddings only through their directions. Both objectives are optimised with Adam and weight decay on the embedding rows.

A trained model is evaluated under four scoring rules:
\begin{align}
  s_{\mathrm{raw}}(u,i)   &= \langle\mathbf{p}_u,\mathbf{q}_i\rangle,
  &
  s_{\cos}(u,i)           &= \cos(\mathbf{p}_u,\mathbf{q}_i),
  \label{eq:scores-plain}\\
  s_{\mathrm{raw}+}(u,i)  &= \langle\mathbf{p}_u,\mathbf{q}_i\rangle+\beta\,\ell_i,
  &
  s_{\mathrm{cr}}(u,i)    &= \cos(\mathbf{p}_u,\mathbf{q}_i)+\beta\,\ell_i.
  \label{eq:scores-prior}
\end{align}
We call $s_{\mathrm{cr}}$ the \emph{count-restored score}. Its form is that of Cheng~\cite[Eq.~53]{Cheng2026}, applied here after training to a shared table. The coefficient $\beta$ is chosen from the grid $\{0,0.01,0.02,0.05,0.1,0.2,0.5,1,2,5\}$ by validation Recall@20, with ties going to the smallest $\beta$. The test split is scored once, at the selected $\beta$. Storing the prior costs one scalar per item. The scoring rule for each objective is fixed in advance; \S\ref{sec:method-intervention} gives the rule and its rationale.

\subsection{Problem Statement and Diagnostic Quantities}
\label{sec:prelim-problem}

We measure whether magnitude follows items or buckets with two rank correlations, computed over items and over non-empty buckets respectively:
\begin{equation}
  \rho_{\mathrm{item}}=\operatorname{Spearman}_{i}\bigl(\|\mathbf{q}_i\|,\,c_i\bigr),
  \qquad
  \rho_{\mathrm{bucket}}=\operatorname{Spearman}_{b}\bigl(\|\mathbf{E}_I[b]\|,\,C_b\bigr).
  \label{eq:rho}
\end{equation}
For an item-level target $t$, let $\mu_b(t)$ be the mean of $t$ over $\mathcal{I}_b$. The \emph{correlation ratio}
\begin{equation}
  \eta^2(t)=\frac{\operatorname{Var}_i\bigl[\mu_{b(i)}(t)\bigr]}{\operatorname{Var}_i\bigl[t_i\bigr]}
  \label{eq:eta}
\end{equation}
is the share of the variance of $t$ that lies between buckets. Variances are taken over items with uniform weight. We use $t=\ell$ for statements about variance and $t=\operatorname{rank}(c)$ for statements about Spearman correlation. Writing $\ell_i=\mu_{b(i)}(\ell)+\varepsilon_i$ splits log-popularity into a bucket-level part and an item-specific residual $\varepsilon_i$.

\begin{table}[t]
  \caption{The problem: raw test Recall@20 (\%) of MF. DirectAU leads BPR on full tables but loses its advantage under GraphHash. $\Delta$ is DirectAU minus BPR in percentage points.}
  \label{tab:problem}
  \small
  \begin{tabular}{@{}lrrrrrr@{}}
    \toprule
    & \multicolumn{3}{c}{Full table} & \multicolumn{3}{c}{GraphHash} \\
    \cmidrule(lr){2-4}\cmidrule(l){5-7}
    Dataset & BPR & DAU & $\Delta$ & BPR & DAU & $\Delta$ \\
    \midrule
    Gowalla    & 16.55 & 18.02 & $+1.47$ & 9.37 & 7.55 & $\mathbf{-1.82}$ \\
    Yelp2018   &  8.02 & 11.10 & $+3.08$ & 3.78 & 2.41 & $\mathbf{-1.37}$ \\
    AmazonBook &  8.38 &  9.57 & $+1.19$ & 4.10 & 4.63 & $+0.53$ \\
    \bottomrule
  \end{tabular}
\end{table}

Table~\ref{tab:problem} states the problem. With full tables, DirectAU outperforms BPR on all three datasets. Under GraphHash, the order reverses on Gowalla and Yelp2018, and on AmazonBook DirectAU's lead shrinks from $1.19$ to $0.53$ points. With LightGCN, BPR already leads on the full Gowalla table ($18.48\%$ vs.\ $18.02\%$), so the premise that DirectAU wins on full tables is specific to MF. Given a fixed sharing map, fixed training protocols and a fixed evaluation protocol, we ask:
\begin{enumerate}
  \item \emph{what} item-level information a shared table cannot represent;
  \item \emph{whether} losing that information accounts for DirectAU's loss of advantage;
  \item \emph{what} restoring it costs.
\end{enumerate}
Section~\ref{sec:method} develops the analysis and the tests that answer these questions.

% ------------------------------------------------------------
% Labels referenced here that must exist elsewhere:
%   sec:method             (Section 4, Diagnostic Methodology)
%   sec:method-intervention (Section 4.6, the intervention as a test)
% Bib keys used: LightGCN, GraphHash, Zhang2020freq, BACO, BPR, DirectAU, Cheng2026
% ------------------------------------------------------------
